\section{四种记忆类型}

该领域已收敛到四种不同的记忆类型，可以类比为大脑的四个不同区域，各自为特定任务而设计。

\subsection{In-context Memory（上下文内记忆）}

Context Window 是 Agent 的\textbf{工作台}——台面上的一切都可以即时访问，模型可以在单次 forward pass 中对其进行推理，无需检索步骤。

但工作台有大小限制：每个 token 都消耗算力和金钱，session 结束后台面会被清空。

上下文中通常包含以下内容：
\begin{itemize}[leftmargin=2em]
    \item \textbf{System Prompt}：Agent 的人设、规则、能力描述、当前日期/用户信息
    \item \textbf{Conversation History}：本次 session 中的对话往来
    \item \textbf{Tool Call Results}：Agent 刚刚调用工具的输出
    \item \textbf{Retrieved Memories}：从外部存储拉取的记忆片段
    \item \textbf{Scratchpad}：中间推理过程（Chain-of-Thought 输出）
\end{itemize}

\begin{warningbox}{Sliding Window 问题}
在长对话中，历史消息不断累积，最终会溢出 context 上限。简单截断最早消息会丢失重要的早期上下文。应采取更好的策略：
\begin{itemize}[leftmargin=1.5em]
    \item \textbf{Summarization}：定期将旧对话压缩为简短摘要
    \item \textbf{Selective Retention}：保留包含关键事实、决策或工具结果的对话轮次，丢弃闲聊
    \item \textbf{Offload to External Memory}：将重要事实提取到 Vector Store，按需检索
\end{itemize}
\end{warningbox}

\subsection{External Memory（外部记忆）}

External Memory 是持久化在模型之外的一切——数据库、向量存储、键值存储和文件。它能跨 session 存活，如果存储得当，Agent 可以回忆起六个月前的事情。

外部存储有两种风格：

\begin{itemize}[leftmargin=2em]
    \item \textbf{Structured Store（精确查找）}：PostgreSQL、Redis、SQLite。通过 key、ID 或 SQL 查询。快速、可预测，适合用户画像、偏好和结构化数据。
    \item \textbf{Vector Store（语义搜索）}：Pinecone、Chroma、pgvector。通过\textbf{语义}查询——"找到与这个概念相似的记忆"。对非结构化笔记和 Episodic Recall 至关重要。
\end{itemize}

\begin{importantbox}{检索设计是核心}
检索步骤是记忆系统的瓶颈。如果你没有检索到正确的记忆，Agent 就像这些记忆根本不存在一样。好的记忆架构是\textbf{20\% 存储 + 80\% 检索设计}。
\end{importantbox}

\subsection{Episodic Memory（情景记忆）}

Episodic Memory 是最被低估的记忆类型。外部记忆存储\textbf{事实}，而情景记忆存储\textbf{事件}——特别是过去行动的结果。

最简单的形式是一个结构化日志：每当 Agent 完成一个任务，它就记录发生了什么。随着时间推移，这个日志变成了一个丰富的自我认知来源，Agent 可以在做决策前查阅它。

一个 Episode 的典型结构：

\begin{itemize}[leftmargin=2em]
    \item \textbf{Task}：任务描述
    \item \textbf{Strategy}：采用的策略
    \item \textbf{Outcome}：执行结果（成功/失败）
    \item \textbf{Lesson}：学到的经验教训
\end{itemize}

\begin{knowledgebox}{Reflection Loop（反思循环）}
当新任务到来时，Agent 检索语义上最相似的历史 Episode，用它们来选择策略。这本质上是从\textbf{个人历史}中进行 Few-shot Learning，而不是从手工构造的数据集。
\end{knowledgebox}

\subsection{Semantic/Parametric Memory（语义/参数记忆）}

这是模型"与生俱来"的记忆——在训练过程中编码到权重里的一切：世界知识、语言模式、推理策略、编码规范和文化常识。

它始终在那里，Agent 永远不需要检索它。但它有严格的限制：

\begin{itemize}[leftmargin=2em]
    \item \textbf{冻结在训练时刻}：模型不知道 cutoff date 之后发生的事
    \item \textbf{运行时无法更新}：不重新训练或微调就无法注入新的永久事实
    \item \textbf{不透明}：你无法精确检查模型"知道"什么或不知道什么
    \item \textbf{容易产生幻觉}：模型用看似合理但实际错误的内容填补知识空白
\end{itemize}

\begin{warningbox}{不要依赖参数记忆处理敏感信息}
对于时间敏感、领域特定或私有的信息，不要依赖 Parametric Memory，应使用外部检索。Parametric Memory 只是在没有更好来源时，作为通用世界知识的兜底方案。
\end{warningbox}

\begin{importantbox}{正确的思维模型}
Parametric Memory 是 Agent 的\textbf{通识教育}；External、Episodic 和 In-context Memory 是 Agent 的\textbf{在职经验}。最好的 Agent 将两者结合使用。
\end{importantbox}

\subsection{本章小结}

四种记忆类型各司其职：In-context Memory 提供即时推理能力，External Memory 提供跨 session 持久化，Episodic Memory 记录行动结果以支持学习，Parametric Memory 提供通用基础知识。

